\section{Hardware Design}
\label{sec:hardware_design}

\subsection{Power Supply Architecture}
\label{sec:power_supply}

\subsubsection{\SI{3.3}{\volt} Rail — Converter Selection}
\label{sec:3v3_converter}

Two hardware revisions are defined.
Revision~1 targets the XIAO ESP32-C3 module and employs the LMR54406 synchronous buck converter at \SI{5}{\volt}, followed by the module's onboard \ac{LDO} (ME6217, \SI{700}{\milli\ampere}) for the final \SI{3.3}{\volt} regulation.
Revision~2 targets the ESP32-S3-MINI-1-N8 on a custom \ac{PCB} and replaces both stages with a single LMR54410 converter delivering \SI{3.3}{\volt} directly.

\begin{table}[H]
  \caption{DC-DC converter parameters — Revision~1 (LMR54406, \SI{5}{\volt} output).}
  \label{tab:converter_rev1}
  \centering
  \begin{tabularx}{\textwidth}{|X|X|}
    \hline
    \textbf{Parameter} & \textbf{Value} \\ \hline \hline
    Topology            & Synchronous Buck \\ \hline
    Device              & LMR54406 (Texas Instruments) \\ \hline
    Input voltage range & \SI{4}{\volt}--\SI{36}{\volt} \\ \hline
    Output voltage      & \SI{5}{\volt} \\ \hline
    Maximum output current & \SI{0.6}{\ampere} \\ \hline
    Quiescent current   & \SI{80}{\micro\ampere} (typ.) \\ \hline
    Package             & SOT-23-6 \\ \hline
  \end{tabularx}
\end{table}

\begin{table}[H]
  \caption{DC-DC converter parameters — Revision~2 (LMR54410, \SI{3.3}{\volt} output).}
  \label{tab:converter_rev2}
  \centering
  \begin{tabularx}{\textwidth}{|X|X|}
    \hline
    \textbf{Parameter} & \textbf{Value} \\ \hline \hline
    Topology            & Synchronous Buck \\ \hline
    Device              & LMR54410 (Texas Instruments) \\ \hline
    Input voltage range & \SI{4}{\volt}--\SI{36}{\volt} \\ \hline
    Output voltage      & \SI{3.3}{\volt} (direct) \\ \hline
    Maximum output current & \SI{1}{\ampere} \\ \hline
    Quiescent current   & \SI{80}{\micro\ampere} (typ.) \\ \hline
    Package             & SOT-23-6 (pin-compatible with LMR54406) \\ \hline
  \end{tabularx}
\end{table}

The LMR54410 feeds \SI{3.3}{\volt} directly to the ESP32-S3-MINI-1-N8 module.
The XIAO board and its onboard \ac{LDO} are eliminated; no topology change is required relative to Revision~1 — only the feedback resistors and current rating differ.

\subsubsection{Input Voltage Source}
\label{sec:input_voltage}

\begin{table}[H]
  \caption{Input voltage source analysis.}
  \label{tab:input_voltage}
  \centering
  \begin{tabularx}{\textwidth}{|X|X|X|X|}
    \hline
    \textbf{Source} & \textbf{Nominal $V_{IN}$} & \textbf{Result} & \textbf{Notes} \\ \hline \hline
    \SI{24}{\volt} industrial supply & \SI{24}{\volt} & OK      & Sole power source for the board. \\ \hline
    \ac{USB} \SI{5}{\volt} (VBUS)    & N/A             & Ignored & VBUS pin is strictly \ac{NC}. \\ \hline
  \end{tabularx}
\end{table}

The system is self-powered from the \SI{24}{\volt} industrial supply.
The \ac{USB} interface is utilised exclusively for data transfer.
The \SI{24}{\volt} supply must be active for the ESP32-S3 to enumerate on a \ac{USB} host.

\subsubsection{Input Protection — \SI{24}{\volt} Rail}
\label{sec:input_protection}

The \SI{24}{\volt} connector is keyed, preventing reverse-polarity insertion.
Industrial \SI{24}{\volt} systems frequently generate inductive voltage spikes exceeding the \SI{36}{\volt} absolute maximum rating of the LMR54410.
A \ac{TVS} diode and a bulk capacitor are mandatory at the input terminal.

\begin{figure}[H]
  \centering
  \begin{circuitikz}[european]
    \draw (0,2) node[left] {24V\_IN}
          to[short, -*] (2,2)
          to[short, -*] (3.5,2)
          to[short, ->] (5.5,2) node[right] {VIN to DC-DC Converters};
    \draw (2,0) to[zD, l=TVS1] (2,2);
    \draw (3.5,2) to[C, l=$C_{bulk}$] (3.5,0);
    \draw (2,0) to[short] (3.5,0);
    \draw (2.75,0) node[sground] {};
  \end{circuitikz}
  \caption{Input protection network on the \SI{24}{\volt} rail.}
  \label{fig:input_protection}
\end{figure}

A unidirectional \ac{TVS} diode (TVS1, SMAJ28A or SMF28A) clamps inductive transients, featuring a standoff voltage of \SI{28}{\volt} and a clamping voltage of approximately \SI{45}{\volt}.

If reverse-polarity protection is required in addition to the keyed connector, a Schottky diode
(e.g.\ SS24, \SI{2}{\ampere}/\SI{40}{\volt}) shall be placed in series on the \texttt{24V\_IN} line
immediately before TVS1.

\subsubsection{\ac{USB}-C Connector Minimum Circuit}
\label{sec:usbc_circuit}

The VBUS line is physically disconnected from the board power network.
GND remains connected to provide a common reference for the D$+$/D$-$ data signals.
Two \SI{5.1}{\kilo\ohm} pull-down resistors on CC1 and CC2 are mandatory; without them a \ac{USB}-C host will not activate the data lines.
CC1 and CC2 must use separate resistors and must not be shorted together.

The \ac{TVS} diode array (e.g.\ USBLC6-2SC6) must connect to the VBUS pin on the \ac{USB} connector
even when VBUS is not used for board power, to provide the necessary reference for the \ac{ESD}
steering diodes.

\begin{figure}[H]
  \centering
  \begin{circuitikz}[european, label distance=1mm]
    \draw (0, 4) node[left] {VBUS} to[short, o-*] (4, 4) to[short, ->] (7.5, 4) node[right] {NC / Floating};

    \draw (0, 2.5) node[left] {CC1} to[R, l=\SI{5.1}{\kilo\ohm}, o-] (2, 2.5) node[sground] {};
    \draw (0, 1.5) node[left] {CC2} to[R, l=\SI{5.1}{\kilo\ohm}, o-] (2, 1.5) node[sground] {};

    \draw (2.5, -1.25) rectangle (5.5, 1.25);
    \node at (4, 0) {\small USBLC6-2SC6};
    \node[right] at (2.5, 0.5) {\tiny 1};
    \node[right] at (2.5, -0.5) {\tiny 3};
    \node[left] at (5.5, 0.5) {\tiny 6};
    \node[left] at (5.5, -0.5) {\tiny 4};
    \node[below] at (4, 1.25) {\tiny 5};
    \node[above] at (4, -1.25) {\tiny 2};

    \draw (4, 4) to[short] (4, 1.25);

    \draw (0, 0.5) node[left] {D$+$} to[short, o-] (2.5, 0.5);
    \draw (5.5, 0.5) to[short, ->] (7.5, 0.5) node[right] {GPIO20 (USB D$+$)};

    \draw (0, -0.5) node[left] {D$-$} to[short, o-] (2.5, -0.5);
    \draw (5.5, -0.5) to[short, ->] (7.5, -0.5) node[right] {GPIO19 (USB D$-$)};

    \draw (0, -1.5) node[left] {GND} to[short, o-] (4, -1.5);
    \draw (4, -1.25) to[short, -*] (4, -1.5);
    \draw (4, -1.5) node[sground] {};

    \draw (0, -2.5) node[left] {Shield} to[R, l=\SI{1}{\mega\ohm}, o-] (4, -2.5) node[sground] {};
    \node[right] at (4.5, -2.5) {or directly to GND};
  \end{circuitikz}
  \caption{\ac{USB}-C connector minimum circuit.}
  \label{fig:usbc_circuit}
\end{figure}

\subsubsection{Current Budget and Power Chain Analysis}
\label{sec:current_budget}

\paragraph{Revision~1 — LMR54406 $\rightarrow$ \SI{5}{\volt} $\rightarrow$ \ac{LDO} $\rightarrow$ \SI{3.3}{\volt}}

The ME6217 is a linear \ac{LDO}; its input current is approximately equal to its output current.
The \SI{5}{\volt} rail must therefore supply the full \SI{3.3}{\volt} load current.

\begin{table}[H]
  \caption{Current budget — Revision~1 (\SI{5}{\volt} rail).}
  \label{tab:current_rev1}
  \centering
  \begin{tabularx}{\textwidth}{|X|X|X|}
    \hline
    \textbf{Consumer} & \textbf{Typical} & \textbf{Peak} \\ \hline \hline
    ESP32-C3 via XIAO (\ac{WiFi} active) & \SI{240}{\milli\ampere} & \SI{500}{\milli\ampere} \\ \hline
    128$\cdot$32 \ac{OLED} (SSD1306)  & \SI{15}{\milli\ampere}--\SI{25}{\milli\ampere} & \SI{30}{\milli\ampere} \\ \hline
    \ac{RS485} driver (MAX485)          & \SI{1}{\milli\ampere}--\SI{5}{\milli\ampere}   & \SI{20}{\milli\ampere} \\ \hline
    \hline
    \textbf{Total on \SI{5}{\volt} rail} & $\approx$\SI{265}{\milli\ampere} & $\approx$\SI{550}{\milli\ampere} \\ \hline
  \end{tabularx}
\end{table}

The LMR54406 is rated at \SI{600}{\milli\ampere}, leaving a peak margin of approximately
\SI{50}{\milli\ampere} — functional but tight. \ac{LDO} thermal dissipation at peak load:

\begin{equation}
  P_{\mathrm{LDO}} = (V_{in} - V_{out}) \cdot I_{load}
                   = (\SI{5}{\volt} - \SI{3.3}{\volt}) \cdot \SI{0.55}{\ampere}
                   = \SI{0.935}{\watt}
  \label{eq:ldo_dissipation}
\end{equation}

At typical load, dissipation is approximately \SI{0.45}{\watt}.
Junction temperature in enclosed enclosures shall be monitored accordingly.

\paragraph{Revision~2 — LMR54410 $\rightarrow$ \SI{3.3}{\volt} direct}

\begin{table}[H]
  \caption{Current budget — Revision~2 (\SI{3.3}{\volt} rail).}
  \label{tab:current_rev2}
  \centering
  \begin{tabularx}{\textwidth}{|X|X|X|}
    \hline
    \textbf{Consumer} & \textbf{Typical} & \textbf{Peak} \\ \hline \hline
    ESP32-S3-MINI-1-N8 (\ac{WiFi} active) & \SI{240}{\milli\ampere} & \SI{500}{\milli\ampere} \\ \hline
    128$\cdot$32 \ac{OLED} (SSD1306)  & \SI{15}{\milli\ampere}--\SI{25}{\milli\ampere} & \SI{30}{\milli\ampere} \\ \hline
    \ac{RS485} driver (MAX3485)         & \SI{1}{\milli\ampere}--\SI{5}{\milli\ampere}   & \SI{20}{\milli\ampere} \\ \hline
    \hline
    \textbf{Total on \SI{3.3}{\volt} rail} & $\approx$\SI{265}{\milli\ampere} & $\approx$\SI{550}{\milli\ampere} \\ \hline
  \end{tabularx}
\end{table}

The LMR54410 is rated at \SI{1000}{\milli\ampere}, yielding a peak margin of \SI{450}{\milli\ampere}.
Elimination of the \ac{LDO} stage removes the dissipation of Equation~\eqref{eq:ldo_dissipation}.
Buck converter efficiency of \SI{80}{\percent}--\SI{90}{\percent} results in substantially lower heat generation compared to Revision~1.

\subsubsection{Output Voltage Setting}
\label{sec:output_voltage}

The feedback reference voltage of both converters is $V_{ref} = \SI{0.8}{\volt}$.
The output voltage is set by the external resistor divider:

\begin{equation}
  V_{out} = V_{ref} \cdot \left(1 + \frac{R_{fbt}}{R_{fbb}}\right)
  \label{eq:vout}
\end{equation}

\paragraph{Revision~1 — \SI{5}{\volt} output (LMR54406)}

\begin{table}[H]
  \caption{Feedback resistors — Revision~1 (\SI{5}{\volt} output).}
  \label{tab:fb_rev1}
  \centering
  \begin{tabularx}{\textwidth}{|X|X|X|}
    \hline
    \textbf{Resistor} & \textbf{Value} & \textbf{Notes} \\ \hline \hline
    $R_{fbt}$ (top, FB to $V_{out}$) & \SI{115.0}{\kilo\ohm} & E96, \SI{1}{\percent} \\ \hline
    $R_{fbb}$ (bottom, FB to GND)    & \SI{22.1}{\kilo\ohm}  & E96, \SI{1}{\percent} \\ \hline
  \end{tabularx}
\end{table}

\begin{equation}
  V_{out} = \SI{0.8}{\volt} \cdot \left(1 + \frac{\SI{115}{\kilo\ohm}}{\SI{22.1}{\kilo\ohm}}\right)
          = \SI{0.8}{\volt} \cdot 6.203 = \SI{4.963}{\volt}
  \label{eq:vout_rev1}
\end{equation}

\paragraph{Revision~2 — \SI{3.3}{\volt} output (LMR54410, datasheet recommended values)}

\begin{table}[H]
  \caption{Feedback resistors — Revision~2 (\SI{3.3}{\volt} output).}
  \label{tab:fb_rev2}
  \centering
  \begin{tabularx}{\textwidth}{|X|X|X|}
    \hline
    \textbf{Resistor} & \textbf{Value} & \textbf{Notes} \\ \hline \hline
    $R_{fbt}$ (top, FB to $V_{out}$) & \SI{69.8}{\kilo\ohm} & E96, \SI{1}{\percent} \\ \hline
    $R_{fbb}$ (bottom, FB to GND)    & \SI{22.1}{\kilo\ohm} & E96, \SI{1}{\percent} \\ \hline
  \end{tabularx}
\end{table}

\begin{equation}
  V_{out} = \SI{0.8}{\volt} \cdot \left(1 + \frac{\SI{69.8}{\kilo\ohm}}{\SI{22.1}{\kilo\ohm}}\right)
          = \SI{0.8}{\volt} \cdot 4.158 = \SI{3.327}{\volt}
  \label{eq:vout_rev2}
\end{equation}

The result of \SI{3.327}{\volt} is within the ESP32's supply tolerance of \SI{3.0}{\volt}--\SI{3.6}{\volt}.
The high-impedance divider (approximately \SI{39}{\micro\ampere}) preserves the converter's low quiescent current advantage.
The datasheet-specified inductor for this configuration is \SI{10}{\micro\henry}.

\subsubsection{Power Indicator \ac{LED} — \SI{3.3}{\volt} Rail}
\label{sec:power_led_3v3}

A green \ac{LED} with series resistor is connected directly on the \SI{3.3}{\volt} rail as a passive power indicator.
No \ac{GPIO} is required; the \ac{LED} is permanently active when the rail is energised.
The design target is \SI{5}{\milli\ampere} ($V_f = \SI{2.0}{\volt}$); the final implementation uses \SI{1.0}{\kilo\ohm} to limit current to approximately \SI{1}{\milli\ampere}.

\begin{equation}
  R_{LED} = \frac{V_{supply} - V_f}{I_{LED}}
          = \frac{\SI{3.3}{\volt} - \SI{2.0}{\volt}}{\SI{5}{\milli\ampere}}
          = \SI{260}{\ohm} \rightarrow \SI{270}{\ohm}
  \label{eq:rled_3v3}
\end{equation}

\begin{table}[H]
  \caption{Power indicator \ac{LED} components — \SI{3.3}{\volt} rail.}
  \label{tab:led_3v3}
  \centering
  \begin{tabularx}{\textwidth}{|X|X|X|}
    \hline
    \textbf{Component} & \textbf{Value} & \textbf{Notes} \\ \hline \hline
    R\_led & \SI{1.0}{\kilo\ohm} (E24, \SI{1}{\percent}) & Limits \ac{LED} current to $\approx$\SI{1}{\milli\ampere}. \\ \hline
    LED1   & Green, 0402 or 0603 & $V_f \approx \SI{2.1}{\volt}$ typ.; standard efficiency type only. \\ \hline
  \end{tabularx}
\end{table}

\subsubsection{Layout Guidelines}
\label{sec:layout_notes}

\begin{itemize}
  \item $R_{fbb}$ shall be placed as close as possible to the FB pin to minimise noise pickup on the feedback node.
  \item The switching node (SW) copper loop area shall be minimised.
  \item The input bulk capacitor ($\geq \SI{100}{\micro\farad}$) shall be placed close to the VIN pin.
\end{itemize}

\subsection{\SI{12}{\volt} Fan Supply}
\label{sec:12v_rail}

\subsubsection{Use Case and Converter}
\label{sec:12v_converter}

The \SI{12}{\volt} rail is derived exclusively from the \SI{24}{\volt} industrial input and supplies the cooling fan.
Since \ac{USB} VBUS is disconnected from the system power network, the \SI{12}{\volt} rail and fan operate only when \SI{24}{\volt} is present.
The same LMR54410 device used for the \SI{3.3}{\volt} rail (Section~\ref{sec:3v3_converter}) is employed, maintaining a single part number across both converters.

\begin{table}[H]
  \caption{DC-DC converter parameters — \SI{12}{\volt} fan supply (LMR54410).}
  \label{tab:converter_12v}
  \centering
  \begin{tabularx}{\textwidth}{|X|X|}
    \hline
    \textbf{Parameter} & \textbf{Value} \\ \hline \hline
    Device          & LMR54410 (Texas Instruments) \\ \hline
    Input           & \SI{24}{\volt} \\ \hline
    Output          & \SI{12}{\volt} \\ \hline
    Load current    & \SI{90}{\milli\ampere} (fan) \\ \hline
    Available headroom & \SI{910}{\milli\ampere} spare — within \SI{1}{\ampere} rating \\ \hline
    Duty cycle      & \SI{50}{\percent} ($12/24$ — nominal operating point) \\ \hline
  \end{tabularx}
\end{table}

\subsubsection{Feedback Resistors — \SI{12}{\volt} Output}
\label{sec:12v_feedback}

Applying Equation~\eqref{eq:vout} with $V_{ref} = \SI{0.8}{\volt}$ and $V_{out} = \SI{12}{\volt}$:

\begin{equation}
  \frac{R_{fbt}}{R_{fbb}} = \frac{V_{out}}{V_{ref}} - 1 = \frac{12}{0.8} - 1 = 14
  \label{eq:fb_ratio_12v}
\end{equation}

\begin{table}[H]
  \caption{Feedback resistors — \SI{12}{\volt} output (datasheet recommended values).}
  \label{tab:fb_12v}
  \centering
  \begin{tabularx}{\textwidth}{|X|X|X|}
    \hline
    \textbf{Resistor} & \textbf{Value} & \textbf{Notes} \\ \hline \hline
    $R_{fbt}$ (top, FB to $V_{out}$) & \SI{309}{\kilo\ohm}  & E96, \SI{1}{\percent} \\ \hline
    $R_{fbb}$ (bottom, FB to GND)    & \SI{22.1}{\kilo\ohm} & E96, \SI{1}{\percent} \\ \hline
  \end{tabularx}
\end{table}

\begin{equation}
  V_{out} = \SI{0.8}{\volt} \cdot \left(1 + \frac{\SI{309}{\kilo\ohm}}{\SI{22.1}{\kilo\ohm}}\right)
          = \SI{0.8}{\volt} \cdot 14.981 = \SI{11.985}{\volt} \quad (-\SI{0.1}{\percent})
  \label{eq:vout_12v}
\end{equation}

The datasheet-specified inductor for this configuration is \SI{33}{\micro\henry}.
The higher inductance relative to the \SI{3.3}{\volt} rail (\SI{10}{\micro\henry}) is expected: at \SI{50}{\percent} duty cycle with \SI{24}{\volt} input the converter operates at a larger volt-second product, requiring greater inductance to maintain the same ripple current specification.

\subsubsection{Power Indicator \ac{LED} — \SI{12}{\volt} Rail}
\label{sec:power_led_12v}

\begin{equation}
  R_{LED} = \frac{V_{supply} - V_f}{I_{LED}}
          = \frac{\SI{12}{\volt} - \SI{2}{\volt}}{\SI{5}{\milli\ampere}}
          = \SI{2}{\kilo\ohm}
  \label{eq:rled_12v}
\end{equation}

\begin{table}[H]
  \caption{Power indicator \ac{LED} components — \SI{12}{\volt} rail.}
  \label{tab:led_12v}
  \centering
  \begin{tabularx}{\textwidth}{|X|X|X|}
    \hline
    \textbf{Component} & \textbf{Value} & \textbf{Notes} \\ \hline \hline
    R\_led & \SI{10}{\kilo\ohm} (E24, \SI{1}{\percent}) & Limits \ac{LED} current to $\approx$\SI{1}{\milli\ampere}. \\ \hline
    LED2   & Green, 0402 or 0603 & Same type as LED1. \\ \hline
  \end{tabularx}
\end{table}

Since the \SI{12}{\volt} rail is active only with \SI{24}{\volt} industrial input present, LED2
implicitly indicates that the \SI{24}{\volt} supply is active.

\subsubsection{Four-Wire Fan Control and Monitoring}
\label{sec:fan_control}

Variable speed control and stall detection are implemented via the standard four-wire fan interface.

\begin{table}[H]
  \caption{Four-wire fan connector pin assignment.}
  \label{tab:fan_pins}
  \centering
  \begin{tabularx}{\textwidth}{|X|X|X|X|}
    \hline
    \textbf{Pin} & \textbf{Function} & \textbf{Connection} & \textbf{Notes} \\ \hline \hline
    1 & GND  & GND              & Common ground \\ \hline
    2 & \SI{12}{\volt} & \SI{12}{\volt} rail & Power from LMR54410 \\ \hline
    3 & Tach & \ac{GPIO} (input)  & Open-drain; 2~pulses/revolution \\ \hline
    4 & \ac{PWM} & \ac{GPIO} (output) & \SI{25}{\kilo\hertz} control signal \\ \hline
  \end{tabularx}
\end{table}

Standard four-wire fans accept \SI{3.3}{\volt} logic ($V_{IH} > \SI{2.8}{\volt}$).
A series resistor protects the ESP32 \ac{GPIO} output:

\begin{figure}[H]
  \centering
  \begin{circuitikz}[european]
    \draw (0,0) node[left] {GPIO\_PWM} to[R, l=$R_{pwm}$ \SI{22}{\ohm}, o-o] (4,0) node[right] {FAN\_PWM};
  \end{circuitikz}
  \caption{\ac{PWM} output protection resistor.}
  \label{fig:fan_pwm}
\end{figure}

The tachometer output is open-collector and requires a pull-up to \SI{3.3}{\volt}:

\begin{figure}[H]
  \centering
  \begin{circuitikz}[european]
    \draw (0,0) node[left] {FAN\_TACH} to[short, o-*] (2,0) to[short, ->] (5,0) node[right] {GPIO\_TACH};
    \draw (2,0) to[R, l=$R_{tach}$ \SI{10}{\kilo\ohm}] (2,2) node[above] {3V3};
    \draw (2,0) to[short, -*] (3.5,0) to[C, l=$C_{tach}$ \SI{100}{\pico\farad}] (3.5,-2) node[sground] {};
    \node[right] at (4, -1.5) {(optional noise filter)};
  \end{circuitikz}
  \caption{Tachometer pull-up and noise filter circuit.}
  \label{fig:fan_tach}
\end{figure}

\subsection{ESP32-S3-MINI-1-N8 Supporting Hardware}
\label{sec:mcu_support}

The ESP32-S3-MINI-1-N8 module requires several supporting circuits that are handled internally on
the XIAO board but must be implemented explicitly on a custom \ac{PCB}.

\subsubsection{Power Supply Decoupling}
\label{sec:decoupling}

\begin{table}[H]
  \caption{VCC decoupling capacitors for the ESP32-S3-MINI-1-N8 module.}
  \label{tab:decoupling}
  \centering
  \begin{tabularx}{\textwidth}{|X|X|X|}
    \hline
    \textbf{Component} & \textbf{Value} & \textbf{Placement} \\ \hline \hline
    C\_bulk   & \SI{10}{\micro\farad} (\SI{6.3}{\volt}) & Close to module VCC pin \\ \hline
    C\_bypass & \SI{100}{\nano\farad} (0402)         & As close as possible to VCC pin \\ \hline
  \end{tabularx}
\end{table}

Both capacitors are mandatory.
The ESP32's internal regulators generate high-frequency switching noise; inadequate decoupling causes unstable operation and \ac{WiFi} performance degradation.

\subsubsection{EN (Enable / Reset) Circuit}
\label{sec:reset_circuit}

The EN pin must be held high for normal operation.
An RC network ensures a clean reset voltage slope; a momentary pushbutton discharges the network to trigger a reset.

\begin{figure}[H]
  \centering
  \begin{circuitikz}[european]
    \draw (0,2) node[above] {3V3} to[R, l=$R_{en}$ \SI{10}{\kilo\ohm}, o-*] (0,0);
    \draw (0,0) to[short, ->] (2,0) node[right] {EN (module pin)};
    \draw (0,0) to[C, l=$C_{en}$ \SI{100}{\nano\farad}] (0,-2) node[sground] {};
    \draw (0,0) to[short, *-] (-2,0) to[nopb, l_=SW\_RESET] (-2,-2) node[sground] {};
  \end{circuitikz}
  \caption{EN pin pull-up and reset circuit.}
  \label{fig:reset_circuit}
\end{figure}

\begin{table}[H]
  \caption{EN circuit components.}
  \label{tab:reset_components}
  \centering
  \begin{tabularx}{\textwidth}{|X|X|X|}
    \hline
    \textbf{Component} & \textbf{Value} & \textbf{Purpose} \\ \hline \hline
    R\_en    & \SI{10}{\kilo\ohm} & Pull-up to \SI{3.3}{\volt} \\ \hline
    C\_en    & \SI{100}{\nano\farad} & RC filter; ensures clean reset voltage slope \\ \hline
    SW\_RESET & Momentary NO & Pulls EN low, triggering reset \\ \hline
  \end{tabularx}
\end{table}

\subsubsection{Boot Strapping Pins}
\label{sec:bootstrap}

Strapping pins are sampled at reset to determine boot mode.
Incorrect logic levels prevent firmware upload or cause boot loops.

\begin{table}[H]
  \caption{Strapping pin requirements for the ESP32-S3-MINI-1-N8.}
  \label{tab:strapping}
  \centering
  \begin{tabularx}{\textwidth}{|X|X|X|}
    \hline
    \textbf{Pin} & \textbf{Required Level} & \textbf{Circuit} \\ \hline \hline
    GPIO0 (BOOT)     & HIGH = normal; LOW = download & \SI{10}{\kilo\ohm} pull-up + SW\_BOOT to GND \\ \hline
    GPIO3            & HIGH                          & \SI{10}{\kilo\ohm} pull-up to \SI{3.3}{\volt} \\ \hline
    GPIO45, GPIO46   & LOW preferred                 & \SI{10}{\kilo\ohm} pull-down to GND \\ \hline
  \end{tabularx}
\end{table}

\begin{figure}[H]
  \centering
  \begin{circuitikz}[european]
    \draw (0,2) node[above] {3V3} to[R, l=$R_{boot}$ \SI{10}{\kilo\ohm}, o-*] (0,0);
    \draw (0,0) to[short, ->] (2,0) node[right] {GPIO0 (BOOT)};
    \draw (0,0) to[nopb, l=SW\_BOOT] (0,-2) node[sground] {};
  \end{circuitikz}
  \caption{BOOT strapping pin circuit.}
  \label{fig:boot_circuit}
\end{figure}

The manual programming sequence (without auto-reset circuit) is as follows:

\begin{enumerate}
  \item Press and hold SW\_BOOT (GPIO0 $\rightarrow$ LOW).
  \item Press and release SW\_RESET (EN pulse $\rightarrow$ reset).
  \item Release SW\_BOOT — the module boots into download mode.
  \item Flash firmware via \ac{USB} or \ac{UART}.
  \item Press SW\_RESET to reboot into normal application mode.
\end{enumerate}

With native \ac{USB}, the esptool auto-reset sequence (RTS/DTR) operates automatically over the \ac{USB} connection; manual button operation is only required for \ac{UART} flashing or where auto-reset is not wired.

\subsubsection{\ac{USB} Interface and \ac{ESD} Protection}
\label{sec:usb_programming}

The ESP32-S3-MINI-1-N8 supports native \ac{USB} (no external \ac{USB}-\ac{UART} bridge required).

\begin{table}[H]
  \caption{\ac{USB} \ac{PHY} pin assignment.}
  \label{tab:usb_pins}
  \centering
  \begin{tabularx}{\textwidth}{|X|X|}
    \hline
    \textbf{Pin} & \textbf{Function} \\ \hline \hline
    GPIO19 & USB D$-$ \\ \hline
    GPIO20 & USB D$+$ \\ \hline
  \end{tabularx}
\end{table}

\ac{ESD} protection is mandatory on the D$+$/D$-$ lines.
The ESP32-S3 \ac{GPIO} pins carry no internal \ac{ESD} protection on the \ac{USB} lines; a plug/unplug event or cable discharge can permanently damage the \ac{USB} \ac{PHY}.

\begin{table}[H]
  \caption{\ac{ESD} protection — \ac{TVS} diode array requirements.}
  \label{tab:esd_requirements}
  \centering
  \begin{tabularx}{\textwidth}{|X|X|}
    \hline
    \textbf{Parameter} & \textbf{Requirement} \\ \hline \hline
    Clamping voltage        & $\leq \SI{5.5}{\volt}$ \\ \hline
    Capacitance per line    & $< \SI{2}{\pico\farad}$ \\ \hline
    \ac{ESD} rating         & IEC~61000-4-2, $\pm\SI{8}{\kilo\volt}$ contact \\ \hline
    Lines protected         & D$+$, D$-$ \\ \hline
  \end{tabularx}
\end{table}

\begin{table}[H]
  \caption{Recommended \ac{TVS} diode array devices for \ac{USB} \ac{ESD} protection.}
  \label{tab:tvs_parts}
  \centering
  \begin{tabularx}{\textwidth}{|l|X|l|l|X|}
    \hline
    \textbf{Part} & \textbf{Manufacturer} & \textbf{Package} & \textbf{$C_{per\,line}$} & \textbf{Notes} \\ \hline \hline
    PRTR5V0U2X     & NXP                & SOT-363  & \SI{0.35}{\pico\farad} & Dual-line clamp \\ \hline
    USBLC6-2SC6    & STMicroelectronics & SOT-23-6 & \SI{1.0}{\pico\farad}  & Widely available \\ \hline
    TPD2E2U06DRLR  & Texas Instruments  & SOT-563  & \SI{0.4}{\pico\farad}  & Compact \\ \hline
  \end{tabularx}
\end{table}

The \ac{TVS} array shall be placed as close as possible to the \ac{USB} connector, before signal
traces reach the module.

\subsubsection{\ac{PCB} Antenna Keep-Out Zone}
\label{sec:antenna_keepout}

The onboard \ac{PCB} antenna of the ESP32-S3-MINI-1-N8 requires a copper-free keep-out zone beneath
and around it. Violations degrade \ac{WiFi} range.

\begin{table}[H]
  \caption{\ac{PCB} antenna keep-out requirements.}
  \label{tab:antenna_keepout}
  \centering
  \begin{tabularx}{\textwidth}{|X|X|}
    \hline
    \textbf{Rule} & \textbf{Value} \\ \hline \hline
    No copper under antenna area             & As defined in datasheet \\ \hline
    No ground plane under antenna            & Required \\ \hline
    Metal object clearance from antenna edge & $\geq \SI{15}{\milli\metre}$ \\ \hline
  \end{tabularx}
\end{table}

\subsubsection{RS-485 Interface}
\label{sec:rs485}

A \SI{3.3}{\volt} \ac{RS485} transceiver (MAX3485 or SN65HVD72) is employed.
The \SI{5}{\volt} MAX485 is not compatible with the \SI{3.3}{\volt} logic levels of the ESP32-S3 and shall not be used.

\begin{table}[H]
  \caption{\ac{RS485} signal and \ac{GPIO} pin assignment.}
  \label{tab:rs485_pins}
  \centering
  \begin{tabularx}{\textwidth}{|l|l|X|}
    \hline
    \textbf{Signal} & \textbf{\ac{GPIO} Pin} & \textbf{Function} \\ \hline \hline
    RO (Receiver Output)    & GPIO5 & \ac{UART} RX \\ \hline
    DI (Driver Input)       & GPIO4 & \ac{UART} TX \\ \hline
    DE/\={RE} (Driver Enable) & GPIO6 & Direction control: HIGH = transmit, LOW = receive \\ \hline
  \end{tabularx}
\end{table}

\begin{figure}[H]
  \centering
  \begin{circuitikz}[european]
    \draw (3, -2) rectangle (5.5, 2);
    \node at (4.25, 1.75) {\small MAX3485};
    \node[right] at (3, 1) {DI};
    \node[right] at (3, 0) {RO};
    \node[right] at (3, -0.75) {DE};
    \node[right] at (3, -1.25) {$\overline{\text{RE}}$};
    \node[left]  at (5.5, 0.5) {A};
    \node[left]  at (5.5, -0.5) {B};
    \draw (0, 1) node[left] {GPIO4 (TX)} to[short, o-] (1.5, 1) to[short, ->] (3, 1);
    \draw (0, 0) node[left] {GPIO5 (RX)} to[short, o-] (1.5, 0) to[short, <-] (3, 0);
    \draw (0, -1) node[left] {GPIO6 (DE)} to[short, o-*] (2, -1);
    \draw (2, -1) |- (2.5, -0.75) to[short, ->] (3, -0.75);
    \draw (2, -1) |- (2.5, -1.25) to[short, ->] (3, -1.25);
    \draw (5.5, 0.5) to[short, -o] (7.5, 0.5) node[right] {D$+$ (Connector)};
    \draw (5.5, -0.5) to[short, -o] (7.5, -0.5) node[right] {D$-$ (Connector)};
  \end{circuitikz}
  \caption{RS-485 transceiver connection diagram.}
  \label{fig:rs485_circuit}
\end{figure}

\subsection{Temperature Sensing}
\label{sec:temp_sensing}

\subsubsection{Sensor Selection}
\label{sec:temp_sensor}

The DS18B20 digital temperature sensor provides remote temperature sensing on the front \ac{PCB} power section.
This sensor communicates over the \ac{OWI} protocol, which supports extended cable runs and reduces wiring to a single data line.

\begin{table}[H]
  \caption{DS18B20 temperature sensor electrical characteristics.}
  \label{tab:ds18b20_specs}
  \centering
  \begin{tabularx}{\textwidth}{|X|X|}
    \hline
    \textbf{Parameter} & \textbf{Value} \\ \hline \hline
    Interface    & \ac{OWI} (Dallas 1-Wire) \\ \hline
    Accuracy     & $\pm\SI{0.5}{\degreeCelsius}$ (from $-\SI{10}{\degreeCelsius}$ to $+\SI{85}{\degreeCelsius}$) \\ \hline
    Range        & $-\SI{55}{\degreeCelsius}$ to $+\SI{125}{\degreeCelsius}$ \\ \hline
    Resolution   & \SI{9}{\bit} to \SI{12}{\bit} (\SI{0.0625}{\degreeCelsius} steps) \\ \hline
    Supply       & \SI{3.0}{\volt} to \SI{5.5}{\volt} \\ \hline
    Address      & Unique 64-bit ROM ID \\ \hline
    Package      & TO-92, SO-8, or waterproof probe \\ \hline
  \end{tabularx}
\end{table}

\subsubsection{Minimum Circuit}
\label{sec:ds18b20_circuit}

The DS18B20 requires a single pull-up resistor on the data line (DQ) to the power rail (VCC).
While parasitic power mode is supported, the RLCV4 employs a standard three-wire connection for maximum reliability during high-temperature monitoring.

\begin{figure}[H]
  \centering
  \begin{circuitikz}[european]
    \draw (0, 2) node[left] {3V3} to[short, o-*] (2, 2)
          to[R, l=$R_{pu}$ \SI{4.7}{\kilo\ohm}] (2, 0)
          to[short, -o] (4, 0) node[right] {DQ (DS18B20 Pin 2)};
    \draw (0, 0) node[left] {GPIO\_TEMP} to[short, o-*] (2, 0);
    \draw (2, 2) to[short, -o] (4, 2) node[right] {VDD (DS18B20 Pin 3)};
    \draw (0, -1) node[left] {GND} to[short, o-o] (4, -1) node[right] {GND (DS18B20 Pin 1)};
  \end{circuitikz}
  \caption{DS18B20 minimum connection circuit.}
  \label{fig:ds18b20_circuit}
\end{figure}
The DS18B20 sensor (SO-8 or TO-92 package) is located on the front \ac{PCB} (\ac{MCPCB}) near the primary heat source.
The \SI{4.7}{\kilo\ohm} pull-up resistor ($R_{pu}$) is located on the main control \ac{PCB}.

The aluminium substrate of the \ac{MCPCB} spreads heat rapidly; sensor placement near the primary heat source minimises thermal response time.

\subsubsection{\ac{OWI} Cable Length Guidelines}
\label{sec:onewire_cable}

\begin{table}[H]
  \caption{\ac{OWI} cable length guidance for the DS18B20 remote sensor.}
  \label{tab:onewire_cable}
  \centering
  \begin{tabularx}{\textwidth}{|l|X|}
    \hline
    \textbf{Cable length} & \textbf{Expected behaviour} \\ \hline \hline
    $< \SI{5}{\metre}$      & Standard \SI{4.7}{\kilo\ohm} pull-up is sufficient for reliable communication. \\ \hline
    \SI{5}{\metre}--\SI{25}{\metre} & Use a lower value pull-up (\SI{2.2}{\kilo\ohm}) to overcome cable capacitance. \\ \hline
    $> \SI{25}{\metre}$     & Use shielded twisted pair cable; consider an active \ac{OWI} bus driver. \\ \hline
  \end{tabularx}
\end{table}
